\section*{参考文献}
\phantomsection
\addcontentsline{toc}{section}{参考文献}

\begin{enumerate}
\def\labelenumi{(\Roman{enumi})}
\tightlist
\item
  Euclidis Elementa, 5 vol., éd. J. L. Heiberg, Lipsiae (Teubner), 1883-88.
\end{enumerate}

(I bis) T. L. HEATH, The thirteen books of Euclid's elements \ldots, 3 vol., Cambridge, 1908.

\begin{enumerate}
\def\labelenumi{(\Roman{enumi})}
\setcounter{enumi}{1}
\tightlist
\item
  Archimedis Opera Omnia, 3 vol., éd. J. L. Heiberg, 2 \(^{e}\) éd., Leipzig (Teubner), 1913-15.
\end{enumerate}

(II bis) Les Œvres complètes d'Archimède, trad. P. Ver Eecke, Paris-Bruxelles (Desclée-de Brouwer), 1921.

\begin{enumerate}
\def\labelenumi{(\Roman{enumi})}
\setcounter{enumi}{2}
\tightlist
\item
  GALILEO GALILEI, Opere, Ristampa della Edizione Nazionale, 20 vol., Firenze (Barbera), 1929-39.
\end{enumerate}

(III bis) Discorsi e Dimonstrazioni Matematiche intorno à due nuove scienze Attenti alla mecanica \& i movimenti locali, del Signor Galileo Galilei Linceo, Filosofo e Matematico primario del Serenissimo Grand Duca di Toscana. In Leida, appresso gli Elsevirii, MDCXXXVIII.

\begin{enumerate}
\def\labelenumi{(\Roman{enumi})}
\setcounter{enumi}{3}
\item
  J. NEPER, Mirifici logarithmorum canonis constructio, Lyon, 1620.
\item
  J. KEPLER, Neue Stereometrie der Fäser, Ostwald's Klassiker, n° 165, Leipzig (Engelmann), 1908.
\item
  B. CAVALIERI: a) Geometria indivisibilibus continuorum quadam ratione promota, Bononiae, 1635 (2 \(^{e}\) éd., 1653); b) Exercitationes geometricae sex, Bononiae, 1647.
\item
  E. TORRICELLI, Opere, 4 vol., éd. G. Loria et G. Vassura, Faenza (Montanari), 1919.
\end{enumerate}

(VII bis) E. TORRICELLI, in G. LORIA, Bibl. Mat. (III), t. I, 1900, p.~78-79.

\begin{enumerate}
\def\labelenumi{(\Roman{enumi})}
\setcounter{enumi}{7}
\item
  G. DE ROBERVAL, Ouvrages de Mathématique (Mémoires de l'Academie Royale des Sciences, t. III), Amsterdam, 1736: a) Observations sur la composition des mouvements, et sur le moyen de trouver les touchantes des lignes courbes, p.~3-67; b) Epistola ad Torricellium, p.~363-399.
\item
  P. GREGORII A SANCTO VINCENTIO, Opus Geometricum Quadraturae Circuli et Sectionum Coni \ldots, 2 vol., Antverpiae, 1647.
\item
  R. DESCARTES, Œuvres, éd. Ch. Adam et P. Tannery, 11 vol., Paris (L. Cerf), 1897-1909.
\end{enumerate}

(X bis) R. DESCARTES, Geometria, trad. latine de Fr.~van Schooten, 2 \(^{e}\) éd., 2 vol., Amsterdam (Elzevir), 1659-61.

\begin{enumerate}
\def\labelenumi{(\Roman{enumi})}
\setcounter{enumi}{10}
\item
  P. FERMAT, Œuvres, Paris (Gauthier-Villars), 1891-1912: a) De linearum curvarum cum lineis rectis comparatione . . . Auctore M. P. E. A. S., t. I, p.~211-254 (trad. française, ibid., t. III, p.181-215); b) Methodus ad disquirendam maximam et minimam, t. I, p.~133-179 (trad. française, ibid., t. III, p.~121-156); c) De aequationum localium transmutatione et emendatione ad multimodam curvilineorum inter se vel cum rectilineis comparationem, cui annectitur proportionis geometricae in quadrandis infinitis parabolis et hyperbolis usus, t. I, p.~255-288 (trad. française, ibid., t. III, p.~216-240).
\item
  B. PASCAL, Œuvres, 14 vol., éd. Brunschvicg, Paris (Hachette), 1904-14: a) Lettre de A. Dettonville à Monsieur A. D. D. S. en lui envoyant la démonstration à la manière des anciens de l'égalité des lignes spirale et parabolique, t. VIII, p.~249-282; b) Lettre de Monsieur Dettonville à Monsieur de Carcavy, t. VIII, p.~334-384; c) Traité des trilignes rectangles et de leurs onglets, t. IX, p.~3-45: d) Traité des sinus du quart de cercle, t. IX, p.~60-76.
\item
  N. MERCATOR, Logarithmotechnia . . . cui nunc accedit vera quadratura hyperbolae . . . Londini, 1668 (reproduced in F. MASÈRES, Scriptores, Logarithmici . . . , vol.~I, London, 1791, p.~167-196).
\item
  LORD BROUNCKER, The Squaring of the Hyperbola by an infinite series of Rational Numbers, together with its Demonstration by the Right Honourable the Lord Viscount Brouncker, Phil. Trans, v. 3 (1668), p.~645-649 (reproduced in F. MASÈRES, Scriptores, Logarithmici \ldots, vol.~I, London, 1791, p.~213-218).
\item
  J. WALLIS, Opera Mathematica, 3 vol., Oxoniae, 1693-95: a) Arithmetica Infinitorum, t. I, p.~355-478.
\end{enumerate}

(XV bis) J. WALLIS, LOGARITHMOTECHNIA NICOLAI MERCATORIS: discoursed in a letter written by Dr.~J. Wallis to the Lord Viscount Brouncker, Phil. Trans., v. 3 (1668), p.~753-759 (reproduced in F. MASÈRES, Scriptores Logarithmici \ldots, vol.~I, London, 1791, p.~219-226).

\begin{enumerate}
\def\labelenumi{(\Roman{enumi})}
\setcounter{enumi}{15}
\tightlist
\item
  CH. HUYGENS, Œuvres complètes, publiées par la Société Hollandaise des Sciences, 19 vol., La Haye (Martinus Nijhoff), 1888-1937: a) Examen de ``Vera Circuli et Hyperboles Quadratura . . .'', t. VI, p.~228-230 (= Journal des Sçavans, juillet 1668); b) Horologium Oscillatorium, t. XVIII; c) Theoremata de Quadratura Hyperboles, Ellipsis et Circuli . . ., t. XI, p.~271-337; d) De Circuli Magnitudine Inventa . . ., t. XII, p.~91-180.
\end{enumerate}

(XVI bis) Christiani Hugenii, Zulichemii Philosophi, vere magni, Dum viveret Zelemii Toparchae, Opera Mechanica, Geomerica, Astronomia et Miscellanea, 4 tomes en 1 vol., Lugd. Batav., 1751.

\begin{enumerate}
\def\labelenumi{(\Roman{enumi})}
\setcounter{enumi}{16}
\tightlist
\item
  J. GREGORY: a) Vera Circuli et Hyperbolae Quadratura . . . , Pataviae, 1667; b) Geometriae Pars Universalis, Pataviae, 1668; c) Exercitationes Geometricae, London, 1668.
\end{enumerate}

(XVII bis) James Gregory Tercentenary Memorial Volume, containing his correspondence with John Collins and his hitherto unpublished mathematical manuscripts . . . , ed.~H. W. Turnbull, London (Bell and Sons), 1939.

\begin{enumerate}
\def\labelenumi{(\Roman{enumi})}
\setcounter{enumi}{17}
\tightlist
\item
  I. BARROW, Mathematical Works, ed.~W. Whewell, Cambridge, 1860.
\end{enumerate}

(XVIII bis) LECTIONES Geometricae: In quibus (praesertim) GENERALIA Curvarum Linearum SYMPTOMATA DECLARANTUR. Auctore ISAAC BARROW \ldots{} Londini \ldots{} M.DC.LXX (= Mathematical Works, p.~156-320).

\begin{enumerate}
\def\labelenumi{(\Roman{enumi})}
\setcounter{enumi}{18}
\item
  I. NEWTON, Opuscula, t. I, Lausanne-Genève (M. M. Bousquet), 1744: a) De Analysi per aequationes numero terminorum infinitas, p.~3-26; b) Methodus fluxionum et serierum infinitarum, p.~29-199 ( \(1^{st}\) publication 1736); c) De Quadratura curvarum, p.~201-244 ( \(1^{st}\) publication 1706); d) Methodus differentialis, p.~271-282 ( \(1^{st}\) publication 1711); e) Commercium Epistolicum, p.~295-420 ( \(1^{st}\) publication 1712, \(2^{nd}\) ed.~1722).
\item
  I. NEWTON, Philosophiae Naturalis Principia Matematica, London, 1687 (new edition, Glasgow, 1871).
\end{enumerate}

(XX bis) I. NEWTON, Mathematical principles of natural philosophy, and his System of the world, transl. into English by A. Motte in 1729, Univ. of California, 1946.

\begin{enumerate}
\def\labelenumi{(\Roman{enumi})}
\setcounter{enumi}{20}
\item
  G.W. LEIBNIZ, Mathematische Schriften, éd. C. I. Gerhardt, 7 vol., Berlin-Halle (Asher-Schmidt), 1849-63: a) Nova Methodus pro Maximis et Minimis \ldots, t. V, p.~220-226 (= Acta Erud. Lips., 1684); b) De Geometria recondita \ldots, t. V, p.~226-233 (= Acta Erud. Lips., 1686); c) Meditatio nova de natura Anguli contactus \ldots, t. VII, p.~326-328 (= Acta Erud. Lips., 1686); d) Generalia de Natura linearum \ldots, t. VII, p.~331-337 (= Acta Erud. Lips., 1692); e) De linea ex lineis numero infinitis \ldots, t. V, p.~266-269 (= Acta Erud. Lips., 1692); f) Nova Calculi differentialis applicatio \ldots, t. V, p.~301-306 (= Acta Erud. Lips., 1694); g) Specimen novum Analyseos \ldots, t. V, p.~350-361 (= Acta Erud. Lips., 1702); h) Continuatio Analyseos Quadraturarum Rationalium, t. V, p.~361-366 (= Acta Erud. Lips., 1703).
\item
  Der Briefwechsel von Gottfried Wilhelm Leibniz mit Mathematikern, herausg. von C. I. Gerhardt, t. I, Berlin (Mayer und Müller), 1899.
\item
  JAKOB BERNOULLI, Opera, 2 vol., Genève (Cramer-Philibert), 1744.
\item
  JOHANN BERNOUILLI, Opera Omnia, 4 vol., Lausanne-Genève (M. M. Bousquet), 1742.
\end{enumerate}

(XXIV bis) JOHANN BERNOUILLI: Die erste Integralrechnung, Ostwald's Klassiker, n° 194, Leipzig-Berlin, 1914; Die Differentialrechnung, Ostwald's Klassiker, n° 211, Leipzig, 1924;

\begin{enumerate}
\def\labelenumi{(\Roman{enumi})}
\setcounter{enumi}{24}
\item
  L. EULER, Opera Omnia: a) Institutiones calculi differentialis, (1) t. X, Berlin (Teubner), 1913; b) Institutiones calculi integralis, (1), t. XI-XIII, Berlin (Teubner), t. 1, 1913-14; c) Commentationes analyticae \ldots{} (1), t. XXIII, Berlin-Zürich (Teubner-O. Füssli), 1938.
\item
  J. d'ALEMBERT: a) Sur les principes métaphysiques du Calcul infinitésimal (Mélanges de littérature, d'histoire et de philosophie, nouv. éd., t. V, Amsterdam, 1768, p.~207-219); b) Encyclopédie, Paris, 1751-65, articles ``Différentiel'' et ``Limite''.
\item
  J.-L. LAGRANGE, Œuvres, Paris (Gauthier-Villars), 1867-1892: a) Théorie des fonctions analytiques, 2 \(^{e}\) éd., t. IX; b) Leçons sur le calcul des fonctions, 2 \(^{e}\) éd., t. X.
\item
  A.-L. CAUCHY, Résumé des leçons données à l'École Royale Polytechnique sur le calcul infinitésimal, Paris, 1823 (= Œuvres complètes, (2), t. IV, Paris (Gauthier-Villars), 1899).
\item
  B. BOLZANO, Schriften, Bd. I: Funktionenlehre, Prag, 1930.
\item
  P. G. LEJEUNE-DIRICHLET, Werke, t. I, Berlin (G. Reimer), 1889.
\item
  B. RIEMANN, Gesammelte Mathematische Werke, 2 \(^{e}\) éd., Leipzig (Teubner), 1892.
\item
  K. WEIERSTRASS, Mathematische Werke, t. II, Berlin (Mayer und Müller), 1895.
\item
  G. DARBOUX, Mémoire sur les fonctions discontinues, Ann. E. N. S., (2) t. IV (1875), p.~57-112.
\item
  C. JORDAN, Traité d'Analyse, 3 \(^{e}\) éd., 3 vol., Paris (Gauthier-Villars), 1909-15.
\end{enumerate}

